\documentclass[10pt,a4paper]{amsart}
\usepackage{amsmath,amssymb,amsthm,graphicx}
% Retained minimal preamble for the Stacks Project chapter more-algebra.tex (More on Algebra).
% The chapter file pinned-original-source.tex carries only \input{preamble} on its first line, so no
% preamble text lives inside the chapter itself.  The Stacks Project uses one single shared preamble for
% every chapter; the snapshot stages it identically as sources/stacks-chow/preamble.tex and
% sources/stacks-etale-cohomology/preamble.tex.  The block below is copied mechanically, line for line,
% from sources/stacks-chow/preamble.tex lines 14-18 (the Xy-pic setup) and 154-217 (the native theorem
% declarations together with the mathematical macros).  Nothing was added, removed or redefined; the
% wrapper class and the amsmath/amssymb/amsthm/graphicx base packages are supplied by the assembler.
% Chapter: more-algebra.tex -- More on Algebra.
% For commutative diagrams we use Xy-pic
\usepackage[all]{xy}

% We use 2cell for 2-commutative diagrams.
\xyoption{2cell}

% Theorem environments.
%
\theoremstyle{plain}
\newtheorem{theorem}[subsection]{Theorem}
\newtheorem{proposition}[subsection]{Proposition}
\newtheorem{lemma}[subsection]{Lemma}

\theoremstyle{definition}
\newtheorem{definition}[subsection]{Definition}
\newtheorem{example}[subsection]{Example}
\newtheorem{exercise}[subsection]{Exercise}
\newtheorem{situation}[subsection]{Situation}

\theoremstyle{remark}
\newtheorem{remark}[subsection]{Remark}
\newtheorem{remarks}[subsection]{Remarks}

\numberwithin{equation}{subsection}

% Macros
%
\def\lim{\mathop{\mathrm{lim}}\nolimits}
\def\colim{\mathop{\mathrm{colim}}\nolimits}
\def\Spec{\mathop{\mathrm{Spec}}}
\def\Hom{\mathop{\mathrm{Hom}}\nolimits}
\def\Ext{\mathop{\mathrm{Ext}}\nolimits}
\def\SheafHom{\mathop{\mathcal{H}\!\mathit{om}}\nolimits}
\def\SheafExt{\mathop{\mathcal{E}\!\mathit{xt}}\nolimits}
\def\Sch{\mathit{Sch}}
\def\Mor{\mathop{\mathrm{Mor}}\nolimits}
\def\Ob{\mathop{\mathrm{Ob}}\nolimits}
\def\Sh{\mathop{\mathit{Sh}}\nolimits}
\def\NL{\mathop{N\!L}\nolimits}
\def\CH{\mathop{\mathrm{CH}}\nolimits}
\def\proetale{{pro\text{-}\acute{e}tale}}
\def\etale{{\acute{e}tale}}
\def\QCoh{\mathit{QCoh}}
\def\Ker{\mathop{\mathrm{Ker}}}
\def\Im{\mathop{\mathrm{Im}}}
\def\Coker{\mathop{\mathrm{Coker}}}
\def\Coim{\mathop{\mathrm{Coim}}}

% Boxtimes
%
\DeclareMathSymbol{\boxtimes}{\mathbin}{AMSa}{"02}

%
% Macros for moduli stacks/spaces
%
\def\QCohstack{\mathcal{QC}\!\mathit{oh}}
\def\Cohstack{\mathcal{C}\!\mathit{oh}}
\def\Spacesstack{\mathcal{S}\!\mathit{paces}}
\def\Quotfunctor{\mathrm{Quot}}
\def\Hilbfunctor{\mathrm{Hilb}}
\def\Curvesstack{\mathcal{C}\!\mathit{urves}}
\def\Polarizedstack{\mathcal{P}\!\mathit{olarized}}
\def\Complexesstack{\mathcal{C}\!\mathit{omplexes}}
% \Pic is the operator that assigns to X its picard group, usage \Pic(X)
% \Picardstack_{X/B} denotes the Picard stack of X over B
% \Picardfunctor_{X/B} denotes the Picard functor of X over B
\def\Pic{\mathop{\mathrm{Pic}}\nolimits}
\def\Picardstack{\mathcal{P}\!\mathit{ic}}
\def\Picardfunctor{\mathrm{Pic}}
\def\Deformationcategory{\mathcal{D}\!\mathit{ef}}

% --- appended 2026-10-12 (Chief): the project own definitions of the reference/slogan/history
% display environments, copied mechanically from the shared Stacks preamble lines 8-12
% (\usepackage{verbatim} plus the three \newenvironment lines, which the chapter preamble keeps
% verbatim from the project preamble). Nothing added, removed or redefined: these environments
% carry no mathematical content and the project itself typesets them as comments. Without these
% three lines a counted statement whose own text uses them fails with Environment slogan undefined.
\usepackage{verbatim}
\newenvironment{reference}{\comment}{\endcomment}
\newenvironment{slogan}{\comment}{\endcomment}
\newenvironment{history}{\comment}{\endcomment}
\usepackage{cleveref}
\crefname{originallemma}{original lemma}{original lemmas}
\Crefname{originallemma}{Original lemma}{Original lemmas}
\makeatletter
\newlabel{lemma-consequence-Artin-Rees}{{28867}{source}{}{sourcepointer-1}{}}
\newlabel{lemma-consequence-Artin-Rees@cref}{{[originallemma][0][]28867}{[1][0][]source}{}{}{}}
\makeatother
\crefname{originallemma}{original lemma}{original lemmas}
\Crefname{originallemma}{Original lemma}{Original lemmas}
\makeatletter
\newlabel{lemma-hom-systems-ML}{{29060}{source}{}{sourcepointer-2}{}}
\newlabel{lemma-hom-systems-ML@cref}{{[originallemma][0][]29060}{[1][0][]source}{}{}{}}
\makeatother
\crefname{originallemma}{original lemma}{original lemmas}
\Crefname{originallemma}{Original lemma}{Original lemmas}
\makeatletter
\newlabel{lemma-iso}{{29168}{source}{}{sourcepointer-3}{}}
\newlabel{lemma-iso@cref}{{[originallemma][0][]29168}{[1][0][]source}{}{}{}}
\makeatother
\crefname{originalremark}{original remark}{original remarks}
\Crefname{originalremark}{Original remark}{Original remarks}
\makeatletter
\newlabel{remark-weird-systems}{{29141}{source}{}{sourcepointer-4}{}}
\newlabel{remark-weird-systems@cref}{{[originalremark][0][]29141}{[1][0][]source}{}{}{}}
\makeatother
\title{Ext groups before and after reduction modulo the powers of an ideal are pro-isomorphic}
\author{The Stacks Project Authors}
\date{stacks-project:more-algebra.tex:local-snapshot-20261012; GNU FDL 1.2 or later}
\begin{document}
\maketitle
\setcounter{subsection}{1}
\noindent\textit{Editorial note.} All delivered mathematical text is the author's own, copied from the pinned Stacks Project chapter file more-algebra.tex; no statement or proof step was rewritten, reordered or removed.\par\medskip
\noindent\textit{Editorial note.} The chapter file pinned-original-source.tex carries only \textbackslash{}input\{preamble\} on its first line, so no preamble text lives inside the chapter. The retained file sources/stacks-more-algebra/retained-preamble.tex is a mechanical, line-for-line copy of the project's single shared preamble as staged in the snapshot, together with the project's own definitions of the reference/slogan/history display environments, which the project itself typesets as comments; nothing was added, removed or redefined by this card and no package was added.\par\medskip
\noindent\textit{Editorial note.} The source URL is the source-level Stacks Project URL, as required when more than one card is produced from one source; the per-claim tag is carried by claim\_key.\par\medskip
\noindent\textit{Editorial note.} The counted statement span includes the author's own \textbackslash{}begin\{reference\} block recording the email correspondence in which the result was found; that block is delivered verbatim and produces no unresolved citation. The author writes '(precise formulation omitted)' and 'A straightforward diagram chase (omitted) shows that the kernel and cokernel of the dotted arrow are annihilated buy I\textasciicircum{}c'; both are kept verbatim, including the typo 'buy' for 'by'. The four references inside the counted proof span are recorded as non-load-bearing navigation pointers.\par\medskip
\noindent\textit{Allowed standard prerequisites (editorial).}\begin{itemize}
\item The category C of systems with the comparison lemma for maps whose kernel and cokernel are annihilated by a fixed power of I. the target category at statement lines 29208-29209 and invoked at proof lines 29228-29244. Is the category in which the isomorphism is asserted and the tool that promotes bounded-annihilation maps to isomorphisms.
\item The Artin-Rees consequence for the homology of a complex, giving the constant c and the almost inverse maps. invoked at proof lines 29224-29225. Produces the constant c that makes the systems (K(n)) and (K\_n) isomorphic.
\item Mittag-Leffler comparison for the Hom systems before and after reduction. invoked at proof lines 29277-29280. Bounds the kernel and cokernel of the two Hom comparison maps by I\textasciicircum{}c-torsion uniformly in n.
\end{itemize}\medskip
\ifdefined\hypertarget\hypertarget{sourcepointer-1}{}\fi
\noindent\textit{Source reference.} Original lemma 28867: sources/stacks-more-algebra/pinned-original-source.tex, source lines 28867--28867. Same-chapter navigation pointer; the statement it names is standard for the section and the argument that consumes it stays inside the counted proof span.\par\medskip
\ifdefined\hypertarget\hypertarget{sourcepointer-2}{}\fi
\noindent\textit{Source reference.} Original lemma 29060: sources/stacks-more-algebra/pinned-original-source.tex, source lines 29060--29060. Same-chapter navigation pointer; the statement it names is standard for the section and the argument that consumes it stays inside the counted proof span.\par\medskip
\ifdefined\hypertarget\hypertarget{sourcepointer-3}{}\fi
\noindent\textit{Source reference.} Original lemma 29168: sources/stacks-more-algebra/pinned-original-source.tex, source lines 29168--29168. Same-chapter navigation pointer; the statement it names is standard for the section and the argument that consumes it stays inside the counted proof span.\par\medskip
\ifdefined\hypertarget\hypertarget{sourcepointer-4}{}\fi
\noindent\textit{Source reference.} Original remark 29141: sources/stacks-more-algebra/pinned-original-source.tex, source lines 29141--29141. Same-chapter navigation pointer; the statement it names is standard for the section and the argument that consumes it stays inside the counted proof span.\par\medskip
\begin{lemma}
\label{lemma-dejong-kollar-kovacs}
\begin{reference}
Email correspondence between Janos Kollar, Sandor Kovacs, and
Johan de Jong of 23/02/2018.
\end{reference}
Let $I$ be an ideal of the Noetherian ring $A$. Let $M$ and $N$
be finite $A$-modules. Write $A_n = A/I^n$, $M_n = M/I^nM$, and
$N_n = N/I^nN$.
For every $i \geq 0$ the objects
$$
\{\Ext^i_A(M, N)/I^n\Ext^i_A(M, N)\}_{n \geq 1}
\quad\text{and}\quad
\{\Ext^i_{A_n}(M_n, N_n)\}_{n \geq 1}
$$
are isomorphic in the category $\mathcal{C}$ of
Remark \ref{remark-weird-systems}.
\end{lemma}

\begin{proof}
Choose a short exact sequence
$$
0 \to K \to A^{\oplus r} \to M \to 0
$$
and set $K_n = K/I^nK$.
For $n \geq 1$ define $K(n) = \Ker(A_n^{\oplus r} \to M_n)$
so that we have exact sequences
$$
0 \to K(n) \to A_n^{\oplus r} \to M_n \to 0
$$
and surjections $K_n \to K(n)$. In fact, by
Lemma \ref{lemma-consequence-Artin-Rees}
there is a $c \geq 0$ and maps $K(n) \to K_n/I^{n - c}K_n$
which are ``almost inverse''.
Since $I^{n - c}K_n \subset K_n[I^c]$ these maps which witness the fact that
the systems $\{K(n)\}_{n \geq 1}$ and $\{K_n\}_{n \geq 1}$
are isomorphic in $\mathcal{C}$.

\medskip\noindent
We claim the systems
$$
\{\Ext^i_{A_n}(K(n), N_n)\}_{n \geq 1}
\quad\text{and}\quad
\{\Ext^i_{A_n}(K_n, N_n)\}_{n \geq 1}
$$
are isomorphic in the category $\mathcal{C}$. Namely, the surjective maps
$K_n \to K(n)$ have kernels annihilated by $I^c$ and therefore determine maps
$$
\Ext^i_{A_n}(K(n), N_n) \to \Ext^i_{A_n}(K_n, N_n)
$$
whose kernel and cokernel are annihilated by $I^c$. Hence the claim
by Lemma \ref{lemma-iso}.

\medskip\noindent
For $i \geq 2$ we have isomorphisms
$$
\Ext^{i - 1}_A(K, N) = \Ext^i_A(M, N)
\quad\text{and}\quad
\Ext^{i - 1}_{A_n}(K(n), N_n) = \Ext^i_{A_n}(M_n, N_n)
$$
In this way we see that it suffices to prove the lemma
for $i = 0, 1$.

\medskip\noindent
For $i = 0, 1$ we consider the commutative diagram
$$
\xymatrix{
0 \ar[r] &
\Hom(M, N) \ar[r] \ar[dd] &
N^{\oplus r} \ar[r]_-\varphi \ar[dd] &
\Hom(K, N) \ar[r] \ar[d] &
\Ext^1(M, N) \ar[r] &
0 \\
& & &
\Hom(K_n, N_n)
\\
0 \ar[r] &
\Hom(M_n, N_n) \ar[r] &
N_n^{\oplus r} \ar[r] &
\Hom(K(n), N_n) \ar[r] \ar[u] &
\Ext^1(M_n, N_n) \ar[r] &
0
}
$$
By Lemma \ref{lemma-hom-systems-ML} we see that the kernel and cokernel of
$\Hom(M, N)/I^n \Hom(M, N) \to \Hom(M_n, N_n)$ and
$\Hom(K, N)/I^n \Hom(K, N) \to \Hom(K_n, N_n)$ and
are $I^c$-torsion for some $c \geq 0$ independent of $n$.
Above we have seen the cokernel of the injective maps
$\Hom(K(n), N_n) \to \Hom(K_n, N_n)$ are annihilated by $I^c$
after possibly increasing $c$. For such a $c$ we obtain maps
$\delta_n : I^c\Hom(K, N)/I^n\Hom(K, N) \to \Hom(K(n), N_n)$
fitting into the diagram (precise formulation omitted).
The kernel and cokernel of $\delta_n$ are annihilated by
$I^c$ after possibly increasing $c$ since we know that the
same thing is true for $\Hom(K, N)/I^n \Hom(K, N) \to \Hom(K_n, N_n)$
and $\Hom(K(n), N_n) \to \Hom(K_n, N_n)$.
Then we can use commutativity of the solid diagram
$$
\xymatrix{
\varphi^{-1}(I^c\Hom(K, N)) \ar[r]_-\varphi \ar[d] &
I^c\Hom(K, N)/I^n\Hom(K, N) \ar[r] \ar[d]^{\delta_n} &
I^c\Ext^1(M, N)/I^n\Ext^1(M, N) \ar[r] \ar@{..>}[d] & 0 \\
N_n^{\oplus r} \ar[r] &
\Hom(K(n), N_n) \ar[r] &
\Ext^1(M_n, N_n) \ar[r] & 0
}
$$
to define the dotted arrow. A straightforward diagram chase
(omitted) shows that the kernel and cokernel of the
dotted arrow are annihilated buy $I^c$ after possibly
increasing $c$ one final time.
\end{proof}

\end{document}
